\documentclass[11pt]{article}
\usepackage{amsmath, amssymb, geometry}
\geometry{letterpaper, margin=1in}

\title{Homework 2.3: Characterizations of Invertible Matrices}
\author{Tashfeen Omran}
\date{September 2026}

\begin{document}
\maketitle

\section*{Part 1: Comprehensive Introduction, Context, and Prerequisites}

\subsection*{Core Concepts}
The topic of ``Characterizations of Invertible Matrices'' centers around a singular, monumental result in linear algebra: The \textbf{Invertible Matrix Theorem (IMT)}. This theorem acts as a massive dictionary, translating mathematical concepts between different distinct dialects of linear algebra. At its core, the topic explores what it means for a square $n \times n$ matrix $A$ to have an inverse, $A^{-1}$, such that $A A^{-1} = A^{-1} A = I_n$. However, instead of just calculating inverses, we use the IMT to show that a matrix being invertible is logically equivalent to a cascade of geometric, algebraic, and functional properties (e.g., linear independence, spanning spaces, and one-to-one transformations). 

\subsection*{Intuition and Derivation}
The intuition behind the IMT relies on the concept of information loss. A linear transformation $T(x) = Ax$ takes vectors from $\mathbb{R}^n$ and maps them to $\mathbb{R}^n$. If $A$ is invertible, the transformation is perfectly reversible; no information is destroyed. 
If information is not destroyed, two distinct vectors cannot map to the same output (hence, the transformation is \textit{one-to-one}, meaning $Ax=0$ has only the trivial solution). Furthermore, if no information is lost in an $n$-dimensional space, the output must cover the entire space (hence, the transformation is \textit{onto}, meaning the columns of $A$ span $\mathbb{R}^n$). 
The IMT formally links these intuitions: because the matrix is square, preserving dimension (one-to-one) guarantees covering the space (onto), and vice-versa. 

\subsection*{Historical Context and Motivation}
Historically, the study of invertibility originated from the need to solve systems of linear equations simultaneously. Mathematicians like Carl Friedrich Gauss and Wilhelm Jordan developed row reduction techniques in the 19th century to systematically find unique solutions to systems. The broader conceptual framework was motivated by the desire to know \textit{before} performing tedious calculations whether a system of equations was guaranteed to have exactly one unique solution. The IMT encapsulates this pursuit, proving that checking one simple property (like the determinant or pivot count) immediately confirms the existence and uniqueness of solutions.

\subsection*{Key Formulas}
\begin{itemize}
    \item \textbf{Inverse Definition:} $A A^{-1} = A^{-1} A = I_n$
    \item \textbf{$2 \times 2$ Inverse:} If $A = \begin{bmatrix} a & b \\ c & d \end{bmatrix}$, then $A^{-1} = \frac{1}{ad-bc}\begin{bmatrix} d & -b \\ -c & a \end{bmatrix}$.
    \item \textbf{Product Inverse:} $(AB)^{-1} = B^{-1} A^{-1}$
    \item \textbf{Transpose Inverse:} $(A^T)^{-1} = (A^{-1})^T$
\end{itemize}

\subsection*{Prerequisites}
To master this topic, you absolutely need absolute fluency in:
\begin{enumerate}
    \item \textbf{Gaussian Elimination:} Finding pivot positions and reducing matrices to Reduced Row Echelon Form (RREF).
    \item \textbf{Definitions of Span and Linear Independence.}
    \item \textbf{Linear Transformations:} Understanding one-to-one (injectivity) and onto (surjectivity).
    \item \textbf{Matrix Multiplication Properties:} Especially associativity $(AB)C = A(BC)$.
\end{enumerate}

\newpage
\section*{Part 2: Detailed Homework Solutions}

\subsection*{Exercises 1--10: Invertibility Checks}
\textit{Objective: Determine if the matrices are invertible with minimal calculation.}

\textbf{1.} $A = \begin{bmatrix} 5 & 7 \\ -3 & -6 \end{bmatrix}$. 
$\det(A) = (5)(-6) - (7)(-3) = -30 + 21 = -9$. Since $-9 \neq 0$, the matrix is \textbf{invertible}.

\textbf{2.} $A = \begin{bmatrix} -4 & 6 \\ 6 & -9 \end{bmatrix}$. 
$\det(A) = (-4)(-9) - (6)(6) = 36 - 36 = 0$. The determinant is zero, so it is \textbf{not invertible}.

\textbf{3.} $A = \begin{bmatrix} 5 & 0 & 0 \\ -3 & -7 & 0 \\ 8 & 5 & -1 \end{bmatrix}$. 
This is a lower triangular matrix. Its determinant is the product of its diagonal entries: $(5)(-7)(-1) = 35 \neq 0$. It is \textbf{invertible}.

\textbf{4.} $A = \begin{bmatrix} -7 & 0 & 4 \\ 3 & 0 & -1 \\ 2 & 0 & 9 \end{bmatrix}$. 
The second column is entirely zeros. Thus, the columns are linearly dependent (a non-trivial combination yields the zero vector). By the IMT, it is \textbf{not invertible}.

\textbf{5.} $A = \begin{bmatrix} 0 & 4 & -7 \\ 1 & 0 & 5 \\ -5 & 8 & -2 \end{bmatrix}$. 
Expanding the determinant along the first column:
\[ \det(A) = -1( (4)(-2) - (-7)(8) ) - 5( (4)(5) - (-7)(0) ) = -1(-8 + 56) - 5(20) = -48 - 100 = -148 \]
Since $\det(A) \neq 0$, the matrix is \textbf{invertible}.

\textbf{6.} $A = \begin{bmatrix} 1 & -5 & -4 \\ 0 & 3 & 4 \\ -3 & 6 & 0 \end{bmatrix}$. 
Performing row operations: $R_3 \leftarrow R_3 + 3R_1$ yields $\begin{bmatrix} 1 & -5 & -4 \\ 0 & 3 & 4 \\ 0 & -9 & -12 \end{bmatrix}$. Notice that the new $R_3$ is exactly $-3R_2$. The rows are linearly dependent, so there will not be 3 pivots. It is \textbf{not invertible}.

\textbf{7.} $A = \begin{bmatrix} -1 & 0 & 2 & 1 \\ -5 & -3 & 9 & 3 \\ 3 & 0 & 1 & -3 \\ 0 & 3 & 1 & 2 \end{bmatrix}$. 
Using column operations, notice that linear dependence exists among the columns, making the determinant zero. It is \textbf{not invertible}.

\textbf{8.} $A = \begin{bmatrix} 1 & 3 & 7 & 4 \\ 0 & 5 & 9 & 6 \\ 0 & 0 & 2 & 8 \\ 0 & 0 & 0 & 0 \end{bmatrix}$. 
This is an upper triangular matrix with a $0$ on the main diagonal. The determinant is $1 \cdot 5 \cdot 2 \cdot 0 = 0$. It is \textbf{not invertible}.

\textbf{9 \& 10.} (Indicated by [T] for Technology). Using standard matrix software to find the RREF of these large matrices reveals that both matrices row reduce to the identity matrix (they have $4$ and $5$ pivots respectively). By the IMT, both are \textbf{invertible}.

\subsection*{Exercises 11--32: Implications of the Invertible Matrix Theorem}
\textit{Objective: Determine True/False or justify statements based on the IMT.}

\textbf{11. (T/F)} If $Ax=0$ has only the trivial solution, then $A$ is row equivalent to the $n \times n$ identity matrix.
\textbf{True.} By the IMT (statements d and b), a trivial null space implies $A$ has $n$ pivots and reduces to $I_n$.

\textbf{12. (T/F)} If there is an $n \times n$ matrix $D$ such that $AD=I$, then there is also an $n \times n$ matrix $C$ such that $CA=I$.
\textbf{True.} By the IMT (statements k and j), a right inverse implies the existence of a left inverse, and both equal $A^{-1}$.

\textbf{13. (T/F)} If the columns of $A$ span $\mathbb{R}^n$, then the columns are linearly independent.
\textbf{True.} By the IMT (statements h and e), spanning $\mathbb{R}^n$ and linear independence are equivalent for square matrices.

\textbf{14. (T/F)} If the columns of $A$ are linearly independent, then the columns of $A$ span $\mathbb{R}^n$.
\textbf{True.} The converse of Ex 13 is also guaranteed by the IMT.

\textbf{15. (T/F)} If $A$ is an $n \times n$ matrix, then the equation $Ax=b$ has at least one solution for each $b$ in $\mathbb{R}^n$.
\textbf{False.} This is only true if $A$ is an \textit{invertible} matrix. Not all $n \times n$ matrices span $\mathbb{R}^n$.

\textbf{16. (T/F)} If the equation $Ax=b$ has at least one solution for each $b$ in $\mathbb{R}^n$, then the solution is unique for each $b$.
\textbf{True.} If it has a solution for every $b$, it spans $\mathbb{R}^n$. By the IMT, $A$ is invertible, which guarantees uniqueness.

\textbf{17. (T/F)} If the equation $Ax=0$ has a nontrivial solution, then $A$ has fewer than $n$ pivot positions.
\textbf{True.} A nontrivial solution means there are free variables, so there cannot be a pivot in every column. Thus, there are fewer than $n$ pivots.

\textbf{18. (T/F)} If the linear transformation $x \mapsto Ax$ maps $\mathbb{R}^n$ into $\mathbb{R}^n$, then $A$ has $n$ pivot positions.
\textbf{False.} Any $n \times n$ matrix maps $\mathbb{R}^n$ \textit{into} $\mathbb{R}^n$. To have $n$ pivots, the mapping must be \textit{onto} $\mathbb{R}^n$.

\textbf{19. (T/F)} If $A^T$ is not invertible, then $A$ is not invertible.
\textbf{True.} By the IMT (statement l), $A$ is invertible if and only if $A^T$ is invertible.

\textbf{20. (T/F)} If there is a $b$ in $\mathbb{R}^n$ such that the equation $Ax=b$ is inconsistent, then the transformation $x \mapsto Ax$ is not one-to-one.
\textbf{True.} Inconsistency means $A$ does not span $\mathbb{R}^n$. By the IMT, if it is not onto, it cannot be one-to-one.

\textbf{21 \& 22.} When is a square upper/lower triangular matrix invertible?
\textbf{Answer:} When all entries on the main diagonal are non-zero. This guarantees $n$ pivot positions.

\textbf{23.} Can a square matrix with two identical columns be invertible?
\textbf{No.} Identical columns mean $c_1 \mathbf{v}_1 + c_2 \mathbf{v}_2 = 0$ has a nontrivial solution (e.g., $c_1=1, c_2=-1$). The columns are linearly dependent.

\textbf{24.} Is it possible for a $5 \times 5$ matrix to be invertible when its columns do not span $\mathbb{R}^5$?
\textbf{No.} By the IMT, spanning the target space is required for invertibility.

\textbf{25.} If $A$ is invertible, then the columns of $A^{-1}$ are linearly independent. Explain why.
\textbf{Answer:} If $A$ is invertible, then $A^{-1}$ is also an invertible matrix (its inverse is $A$). By the IMT, any invertible matrix has linearly independent columns.

\textbf{26.} If $C$ is $6 \times 6$ and $Cx=v$ is consistent for every $v \in \mathbb{R}^6$, is it possible that for some $v$, the equation $Cx=v$ has more than one solution?
\textbf{No.} Consistency for every $v$ means $C$ spans $\mathbb{R}^6$. By the IMT, $C$ is one-to-one, guaranteeing unique solutions.

\textbf{27.} If the columns of a $7 \times 7$ matrix $D$ are linearly independent, what can you say about solutions of $Dx=b$?
\textbf{Answer:} By the IMT, linear independence means $D$ is invertible, so $Dx=b$ has exactly one unique solution for every $b \in \mathbb{R}^7$.

\textbf{28.} If $n \times n$ matrices $E$ and $F$ have the property that $EF=I$, then $E$ and $F$ commute. Explain why.
\textbf{Answer:} By the IMT, if $EF=I$, then $F$ must be the true inverse of $E$ ($F = E^{-1}$). By definition of an inverse, $FE = E^{-1}E = I$. Thus $EF = FE$.

\textbf{29.} If the equation $Gx=y$ has more than one solution for some $y \in \mathbb{R}^n$, can the columns of $G$ span $\mathbb{R}^n$?
\textbf{No.} Multiple solutions mean $G$ is not one-to-one. By the IMT, a square matrix that is not one-to-one cannot span $\mathbb{R}^n$.

\textbf{30.} If the equation $Hx=c$ is inconsistent for some $c$ in $\mathbb{R}^n$, what can you say about the equation $Hx=0$?
\textbf{Answer:} Inconsistency means $H$ does not span $\mathbb{R}^n$. By the IMT, $H$ is not invertible, which means $Hx=0$ must have non-trivial solutions.

\textbf{31.} If an $n \times n$ matrix $K$ cannot be row reduced to $I_n$, what can you say about the columns of $K$?
\textbf{Answer:} By the IMT, the matrix is not invertible. Therefore, its columns are linearly dependent and do not span $\mathbb{R}^n$.

\textbf{32.} If $L$ is $n \times n$ and $Lx=0$ has the trivial solution, do the columns of $L$ span $\mathbb{R}^n$?
\textbf{Yes.} A trivial null space means $L$ has $n$ pivots. By the IMT, it must span $\mathbb{R}^n$.

\subsection*{Exercises 33--48: Proofs and Linear Transformations}
\textbf{33.} Verify: If $AB=I$, then $A$ and $B$ are invertible, with $B=A^{-1}$ and $A=B^{-1}$.
\textbf{Proof:} Because $A$ and $B$ are square and $AB=I$, statement (k) of the IMT implies that $A$ is invertible. Let $A^{-1}$ be its inverse. Multiplying $AB=I$ on the left by $A^{-1}$ yields $A^{-1}AB = A^{-1}I \implies B = A^{-1}$. Since $A$ is invertible, its inverse $B$ is also invertible, and $B^{-1} = (A^{-1})^{-1} = A$.

\textbf{34.} Why do columns of $A^2$ span $\mathbb{R}^n$ when columns of $A$ are linearly independent?
\textbf{Proof:} If $A$ has independent columns, $A$ is invertible. The product of two invertible matrices ($A \cdot A = A^2$) is invertible. By the IMT, an invertible matrix $A^2$ must span $\mathbb{R}^n$.

\textbf{35 \& 36.} Show that if $AB$ is invertible, so is $A$ and so is $B$.
\textbf{Proof:} If $AB$ is invertible, there exists $W$ such that $(AB)W = I$ and $W(AB) = I$. 
For $A$: $(AB)W = I \implies A(BW) = I$. Since $A$ multiplied by $(BW)$ yields the identity, $A$ is invertible by the IMT.
For $B$: $W(AB) = I \implies (WA)B = I$. Since $(WA)$ multiplied by $B$ yields the identity, $B$ is invertible by the IMT.

\textbf{37 \& 38.} Connections between equations and one-to-one properties.
\textbf{Answer (37):} If $Ax=b$ has multiple solutions, $A$ is not invertible, so the transformation $x \mapsto Ax$ is not one-to-one.
\textbf{Answer (38):} If $x \mapsto Ax$ is one-to-one, $A$ is invertible. Thus, the transformation also maps $\mathbb{R}^n$ \textit{onto} $\mathbb{R}^n$.

\textbf{39 \& 40.} Logic without referencing theorems by name.
\textbf{Answer (39):} If $Ax=b$ has a solution for all $b$, every row has a pivot. Since it is square, every column has a pivot. This means no free variables exist, making the solution unique.
\textbf{Answer (40):} If $Ax=0$ has only the trivial solution, there are no free variables, so every column has a pivot. Because it is square, every row has a pivot, guaranteeing a solution for all $b$.

\textbf{41.} Show $T(x_1, x_2) = (-9x_1 + 7x_2, 4x_1 - 3x_2)$ is invertible and find $T^{-1}$.
\textbf{Solution:} Standard matrix $A = \begin{bmatrix} -9 & 7 \\ 4 & -3 \end{bmatrix}$. 
$\det(A) = 27 - 28 = -1$.
$A^{-1} = \frac{1}{-1} \begin{bmatrix} -3 & -7 \\ -4 & -9 \end{bmatrix} = \begin{bmatrix} 3 & 7 \\ 4 & 9 \end{bmatrix}$. 
Thus, $T^{-1}(x_1, x_2) = (3x_1 + 7x_2, 4x_1 + 9x_2)$.

\textbf{42.} Show $T(x_1, x_2) = (6x_1 - 8x_2, -5x_1 + 7x_2)$ is invertible and find $T^{-1}$.
\textbf{Solution:} Standard matrix $A = \begin{bmatrix} 6 & -8 \\ -5 & 7 \end{bmatrix}$. 
$\det(A) = 42 - 40 = 2$.
$A^{-1} = \frac{1}{2} \begin{bmatrix} 7 & 8 \\ 5 & 6 \end{bmatrix}$. 
Thus, $T^{-1}(x_1, x_2) = (3.5x_1 + 4x_2, 2.5x_1 + 3x_2)$.

\textbf{43.} Explain why an invertible transformation $T$ is one-to-one and onto using equations (1) and (2).
\textbf{Solution:} By (1) $S(T(x)) = x$. If $T(u) = T(v)$, then $S(T(u)) = S(T(v)) \implies u = v$, meaning it is one-to-one. By (2) $T(S(x)) = x$. For any target $y$, let $x = S(y)$. Then $T(x) = y$, proving it is onto.

\textbf{44-48.} These proofs rely on identical algebraic manipulations of transformations. E.g., for Ex 47 (uniqueness of inverse), assume a second inverse $U$ exists. Since $T$ is onto, any $v = T(x)$. Then $S(v) = S(T(x)) = x$ and $U(v) = U(T(x)) = x$. Thus $S(v) = U(v)$, proving uniqueness.

\subsection*{Exercises 49--53: Numerical and Condition Number [T]}
\textit{Objective: Observe the breakdown of theoretical math in floating-point computation.}
\textbf{49.} Solving the systems algebraically yields vastly different results despite extremely similar coefficients. 
System 1 yields $x = (3.94, 0.49)$. System 2 yields $x = (2.9, 2.0)$. A $0.05\%$ change in input yields a $>300\%$ change in output because the matrix is ``ill-conditioned'' (its columns are nearly parallel).
\textbf{50-53.} These computational exercises require the use of MATLAB or Python to calculate the condition number, $\kappa = \|A\| \|A^{-1}\|$. A large condition number (like those found in 5th and 12th order Hilbert matrices) signifies a matrix that is mathematically invertible but practically useless in a computer because infinitesimal roundoff errors explode, destroying precision.

\section*{Part 3: In-Depth Analysis of Problems and Techniques}

\subsection*{A) Problem Types and General Approach}
\begin{itemize}
    \item \textbf{Setup Only / Concept Check (Ex 11-32):} These true/false problems test pure logical deduction. The strategy is to map the given statement to a letter in the IMT, and then follow the equivalence circle to the desired conclusion.
    \item \textbf{Matrix Inversion Problems (Ex 41-42):} These require extracting the standard matrix from a linear transformation description, calculating the determinant, and applying the 2D inverse formula.
    \item \textbf{Theoretical Proofs (Ex 33-36):} These require leveraging the associativity of matrix multiplication $(AB)C = A(BC)$ to isolate variables and show that a matrix meets the definition of invertibility.
\end{itemize}

\subsection*{B) Key Algebraic and Calculus Manipulations}
\begin{itemize}
    \item \textbf{Determinant Expansion:} Used in Ex 5. Expanding along a column with the most zeros minimizes arithmetic. 
    \item \textbf{Direct Row Operation Check:} In Ex 6, stopping row reduction early once it was obvious that $R_3 = -3R_2$ saved time. Recognizing linearly dependent rows prevents the need for full echelon form computation.
    \item \textbf{Matrix Factoring / Associativity:} In Ex 35, rewriting $(AB)W = I$ as $A(BW) = I$ is a crucial trick. Because matrix multiplication is strictly non-commutative, you cannot arbitrarily move terms, but you \textit{can} regroup brackets.
\end{itemize}

\section*{Part 4: ``Cheatsheet'' and Tips for Success}

\subsection*{Summary of Formulas}
\begin{itemize}
    \item $A^{-1} = \frac{1}{\det(A)} \begin{bmatrix} d & -b \\ -c & a \end{bmatrix}$
    \item $(AB)^{-1} = B^{-1} A^{-1}$ (The ``Socks and Shoes'' rule: you put socks on before shoes, but take shoes off before socks).
    \item $(A^T)^{-1} = (A^{-1})^T$
\end{itemize}

\subsection*{Cheats and Shortcuts}
\begin{itemize}
    \item \textbf{Triangular Trick:} If a matrix is upper or lower triangular, do not calculate the determinant normally. Just multiply the entries on the main diagonal. If there is a $0$ anywhere on the diagonal, it is NOT invertible.
    \item \textbf{Identical Row/Column Trick:} If you spot two identical rows or columns (or multiples of each other), the determinant is zero instantly.
\end{itemize}

\subsection*{Common Pitfalls and Traps}
\begin{itemize}
    \item \textbf{Non-Commutativity:} Assuming $AB = BA$. Never assume this unless explicitly told.
    \item \textbf{Square Restriction Trap:} Applying the IMT to non-square ($m \times n$) matrices. The IMT ONLY applies to $n \times n$ matrices. If a matrix is $3 \times 4$, it cannot be invertible, even if its columns span $\mathbb{R}^3$.
\end{itemize}

\section*{Part 5: Conceptual Synthesis and The ``Big Picture''}

\subsection*{A) Thematic Connections}
The core theme of this topic is \textbf{Equivalence and Duality}. In mathematics, finding that two seemingly unrelated concepts are equivalent allows us to solve hard problems using easy tools. The IMT bridges Algebra (solving equations), Geometry (spanning multidimensional spaces), and Calculus/Functions (injectivity and surjectivity). 

\subsection*{B) Forward and Backward Links}
\textbf{Backward:} This chapter is the culmination of everything in Chapter 1. The algorithms of row reduction and the definitions of span and linear independence were built specifically to be united under this theorem.
\textbf{Forward:} The IMT is foundational for \textit{Eigenvalues} (a matrix $A$ has 0 as an eigenvalue if and only if $A$ is not invertible). In Multivariable Calculus, the invertibility of the \textit{Jacobian Matrix} is required to apply the Inverse Function Theorem, which dictates whether a complex coordinate transformation can be reversed.

\section*{Part 6: Real-World Application and Modeling}

\subsection*{A) Concrete Scenarios in Finance and Stochastic Calculus}
\begin{enumerate}
    \item \textbf{Markowitz Portfolio Optimization:} To find the optimal weights of a diversified financial portfolio, quantitative analysts must invert the Covariance Matrix ($\Sigma$) of asset returns. If the matrix is not invertible, it mathematically indicates that one asset is perfectly correlated with a combination of others (a redundant asset).
    \item \textbf{Implicit Finite Difference Methods (Black-Scholes):} When pricing exotic derivative securities using numerical partial differential equations, the time-stepping algorithm requires solving a system of equations backward in time: $V_{t} = A V_{t+1} + b$. The matrix $A$ must be inverted at every single micro-second time step.
\end{enumerate}

\subsection*{B) Model Problem Setup: Markowitz Optimization}
\textbf{Model:} An asset manager is allocating funds among $n$ assets.
\textbf{Variables:} Let $\mathbf{w}$ be the $n \times 1$ vector of portfolio weights. Let $\Sigma$ be the $n \times n$ covariance matrix of historical returns. Let $\mathbf{1}$ be an $n \times 1$ vector of ones.
\textbf{Objective:} Minimize portfolio variance, given by $\mathbf{w}^T \Sigma \mathbf{w}$, subject to the constraint that all weights sum to 100\% ($\mathbf{w}^T \mathbf{1} = 1$).
\textbf{Equation to Solve:} Using Lagrange multipliers, the unscaled optimal weights are found by solving the linear system:
\[ \Sigma \mathbf{u} = \mathbf{1} \]
To find $\mathbf{u}$, the computer must calculate $\mathbf{u} = \Sigma^{-1} \mathbf{1}$. By the IMT, this relies entirely on $\Sigma$ being invertible, which requires the historical asset returns to form a linearly independent set.

\section*{Part 7: Common Variations and Untested Concepts}
\subsection*{Block Matrix Inversion}
While the homework covers $n \times n$ matrices generally, a highly common advanced exam question involves finding the inverse of a matrix partitioned into blocks.
\textbf{Concept:} If a matrix is partitioned as $M = \begin{bmatrix} A & B \\ 0 & D \end{bmatrix}$, where $A$ and $D$ are square invertible matrices, its inverse is given by:
\[ M^{-1} = \begin{bmatrix} A^{-1} & -A^{-1} B D^{-1} \\ 0 & D^{-1} \end{bmatrix} \]
\textbf{Example Problem:} Find the inverse of $M = \begin{bmatrix} 2 & 3 \\ 0 & 4 \end{bmatrix}$ using block inversion.
\textit{Solution:} Here $A = \begin{bmatrix} 2 \end{bmatrix}$ and $D = \begin{bmatrix} 4 \end{bmatrix}$. Thus $A^{-1} = \begin{bmatrix} 1/2 \end{bmatrix}$ and $D^{-1} = \begin{bmatrix} 1/4 \end{bmatrix}$. 
The top right block is $-(1/2)(3)(1/4) = -3/8$.
Thus, $M^{-1} = \begin{bmatrix} 1/2 & -3/8 \\ 0 & 1/4 \end{bmatrix}$.

\section*{Part 8: Advanced Diagnostic Testing: ``Find the Flaw''}
\textit{Instructions for the user: The following five problems contain a complete solution with a subtle, critical error. Your task is to act as the grader, locate the error, explain why it is wrong, and correct it. (An answer key is provided at the bottom).}

\textbf{Problem 1:} Determine if the matrix $A = \begin{bmatrix} 4 & 5 \\ 2 & 2 \end{bmatrix}$ is invertible and find its inverse if it exists.
\newline
\textbf{Flawed Solution:} 
We calculate the determinant. $\det(A) = (4)(2) + (5)(2) = 8 + 10 = 18$. Since the determinant is not zero, the matrix is invertible.
The inverse is calculated using the standard formula:
$A^{-1} = \frac{1}{18}\begin{bmatrix} 2 & -5 \\ -2 & 4 \end{bmatrix}$.

\vspace{0.5cm}
\textbf{Problem 2:} Let $A$ and $B$ be $n \times n$ invertible matrices. Simplify the expression $(A B)^{-1} A$.
\newline
\textbf{Flawed Solution:} 
Distribute the inverse across the product: $(A B)^{-1} A = A^{-1} B^{-1} A$.
Since matrix multiplication is commutative, we can rearrange the terms: $A^{-1} A B^{-1}$.
Since $A^{-1} A = I$, this simplifies to $I B^{-1} = B^{-1}$.

\vspace{0.5cm}
\textbf{Problem 3:} Suppose $A$ is a $3 \times 4$ matrix whose columns span $\mathbb{R}^3$. Is $A$ invertible?
\newline
\textbf{Flawed Solution:} 
Because the columns of $A$ span $\mathbb{R}^3$, the transformation $x \mapsto Ax$ maps onto the target space. By statement (h) of the Invertible Matrix Theorem, any matrix that spans its target space is invertible. Therefore, $A$ is invertible.

\vspace{0.5cm}
\textbf{Problem 4:} Let $T: \mathbb{R}^n \to \mathbb{R}^n$ be a linear transformation. If $T(u) = T(v)$ for some distinct vectors $u$ and $v$, can $T$ map onto $\mathbb{R}^n$?
\newline
\textbf{Flawed Solution:} 
Yes. The condition $T(u) = T(v)$ simply means that two inputs map to the same output, meaning the transformation is not one-to-one. However, a transformation can still be onto even if it is not one-to-one. Therefore, it is entirely possible that $T$ can still map onto $\mathbb{R}^n$.

\vspace{0.5cm}
\textbf{Problem 5:} Prove that if $A^2 = 0$ for an $n \times n$ matrix $A$, then $A$ is not invertible.
\newline
\textbf{Flawed Solution:} 
Assume we have the equation $A^2 = 0$. We can multiply both sides of the equation by $A^{-1}$ to get $A^{-1} A^2 = A^{-1} 0$.
This simplifies to $(A^{-1} A) A = 0$, which means $I A = 0$, so $A = 0$.
Since the zero matrix is explicitly not invertible (its determinant is 0), $A$ is not invertible.

\vspace{1cm}
\hrule
\vspace{0.5cm}
\subsubsection*{Answer Key for Grader}
\textbf{Key 1:} The determinant formula is $ad-bc$, not $ad+bc$. The determinant is $(4)(2) - (5)(2) = 8 - 10 = -2$. The correct inverse is $A^{-1} = -\frac{1}{2}\begin{bmatrix} 2 & -5 \\ -2 & 4 \end{bmatrix}$. \\
\textbf{Key 2:} Matrix multiplication is NOT commutative, and the inverse of a product reverses the order. $(AB)^{-1} = B^{-1}A^{-1}$. Thus $(AB)^{-1} A = B^{-1}A^{-1}A = B^{-1}I = B^{-1}$. The flawed solution got the right answer using two mathematical atrocities that canceled each other out. \\
\textbf{Key 3:} The Invertible Matrix Theorem \textit{only} applies to square ($n \times n$) matrices. A $3 \times 4$ matrix is physically impossible to invert because it maps between spaces of different dimensions. \\
\textbf{Key 4:} Because the transformation maps $\mathbb{R}^n \to \mathbb{R}^n$ (square matrix mapping), the IMT dictates that one-to-one and onto are logically tied together. If it is not one-to-one, it physically cannot be onto. \\
\textbf{Key 5:} Circular logic. The proof starts by multiplying by $A^{-1}$, which assumes $A$ is invertible, in an attempt to prove it is not invertible. A correct proof would use determinants: $\det(A^2) = \det(0) \implies (\det(A))^2 = 0 \implies \det(A) = 0$.

\end{document}